\begin{textsample}{POS Dim 3 – se – Score 12.00 – se/se\_inf\_9.txt}  \label{ex:f3_pos_213}
2.4.2 Os Objetivos de Aprendizagem e Desenvolvimento na Educação \textbf{Infantil} O organizador curricular na Educação \textbf{Infantil} está estruturado em 03 fases com faixas etárias distintas, a saber : - \textbf{Crianças} com idade entre 0 e 1 ano e 6 \textbf{meses} - \textbf{Crianças} com idade entre 1 ano e 7 \textbf{meses} e 3 anos e 11 \textbf{meses} - \textbf{Crianças} com idade entre 4 anos e 5 anos e 11 \textbf{meses} As proposições de objetivos de aprendizagem e desenvolvimento estarão dispostos em quadros por Campos de Experiências e adequados à faixa etária. 2.5 ORGANIZADOR CURRICULAR DA EDUCAÇÃO \textbf{INFANTIL} 1.
GRUPO 1 – \textbf{BEBÊS} ( DE 0 A 1 ANO E 6 \textbf{MESES} ) CAMPO 1 : O EU, O OUTRO E O NÓS Direitos de aprendizagem e desenvolvimento : Conviver, \textbf{Brincar}, Participar, Expressar, Conhecer -se e Explorar.
O campo de experiências O Eu, o outro e o nós permite que as \textbf{crianças} possam compreender a si mesmas e ao outro em qualificadas relações humanas de convivência, contribuindo para construção da identidade, dos valores éticos e da cidadania.
Isto posto, torna -se relevante a interlocução com os demais campos de experiências que as \textbf{crianças} precisam conectar nas interações e \textbf{brincadeiras}.
Os \textbf{bebês}, desde bem \textbf{pequenos}, têm iniciativas de buscar interagir com os \textbf{adultos} e as outras \textbf{crianças}.
As relações de confiança e segurança são essenciais para motivar sua autoiniciativa de interação por meio de seu corpo e sentidos, para que, nos momentos de exploração, aprendam sobre o mundo à sua \textbf{volta}.
Ao serem convidados a \textbf{brincar} próximos a outras \textbf{crianças} ou a interagir com elas ou com seus(as) professores(as), os \textbf{bebês} descobrem diferentes formas de se expressarem e se comunicarem.
Por meio de situações de interação livre ou intencionalmente planejadas e vivenciadas com professores(as) e outros \textbf{adultos}, nos quais confiam, os \textbf{bebês} continuam suas descobertas sobre si mesmos, percebendo -se como um ser individual, com necessidades e \textbf{desejos} próprios ; aprendem a participar e colaborar em situações de convivência, em contato com \textbf{colegas}, podendo ser em dupla, trio ou em grupo.
Valorizando e descobrindo diferentes formas de estar com os outros, os \textbf{bebês} começam a perceber que são capazes de conseguir reações específicas a partir de suas ações, e que suas ações têm efeitos nas outras pessoas.

% matched lemmas: adulto, bebê, brincadeira, brincar, colega, criança, desejo, infantil, mês, pequeno, volta
\end{textsample}
